\documentclass[11pt]{article}
\usepackage{amsmath}
\usepackage{graphicx}
\usepackage{booktabs}
\usepackage{hyperref}

\title{Quantum error correction below the surface code threshold}
\author{Google Quantum AI and Collaborators}
\date{arXiv:2408.13687 reduced Bilin parser fixture}

\newcommand{\logical}{p_{\mathrm{L}}}
\newcommand{\physical}{p_{\mathrm{phys}}}

\begin{document}
\maketitle

\begin{abstract}
This reduced fixture is based on the public arXiv record for
arXiv:2408.13687 and is used only to exercise Bilin parser structure.
\end{abstract}

\section{Below-threshold memory}
\label{sec:memory}

Quantum error correction combines physical qubits into a logical memory. The
fixture keeps the structure of a real arXiv physics paper while using compact
text suitable for deterministic regression tests.

\begin{equation}
\label{eq:suppression}
\logical(d+2) \approx \frac{\logical(d)}{\Lambda}
\end{equation}

\section{Real-time decoding}
\label{sec:decoder}

The paper reports real-time decoding alongside distance scaling. Bilin should
preserve this paragraph as source Markdown and keep citation placeholders such
as Ref.~\cite{acharya2024threshold}.

\begin{figure}
\centering
\includegraphics[width=.72\linewidth]{assets/surface-code-memory.svg}
\caption{A reduced surface-code memory schematic for parser regression.}
\label{fig:memory}
\end{figure}

\begin{table}
\centering
\caption{Representative structured values retained by the fixture.}
\label{tab:values}
\begin{tabular}{lll}
\toprule
Quantity & Value & Role \\
\midrule
Distance & 7 & logical memory \\
Cycles & $10^6$ & real-time decoder run \\
\bottomrule
\end{tabular}
\end{table}

\begin{thebibliography}{1}
\bibitem{acharya2024threshold}
Google Quantum AI and Collaborators.
Quantum error correction below the surface code threshold.
arXiv:2408.13687, 2024.
\end{thebibliography}

\end{document}
